\chapter*{Conclusiones y recomendaciones}
\addcontentsline{toc}{chapter}{Conclusiones y recomendaciones}
\markboth{Conclusiones y recomendaciones}{Conclusiones y recomendaciones}

\section*{Conclusiones}

El presente trabajo cumplió con el objetivo general planteado: se desarrolló \emph{EduFEM}, un software educativo de escritorio para el análisis por el Método de Elementos Finitos (MEF) de problemas bidimensionales de elasticidad lineal en tensión plana y deformación plana, con elementos cuadriláteros isoparamétricos Q4 y Q9. La herramienta integra, en un único entorno coherente, las tres dimensiones que se propusieron como ejes del proyecto: la formulación matemática del método, la construcción y visualización interactiva del modelo, y la verificación numérica de los resultados. El trabajo regresa así a la contradicción que lo originó: la doble brecha entre la abstracción de los textos clásicos del MEF y la opacidad del software profesional como caja negra. EduFEM la atiende ofreciendo un entorno que conserva el rigor numérico —verificado y validado— a la vez que hace observable, sobre el modelo concreto del estudiante, el procesamiento intermedio que ambos extremos dejaban fuera de su alcance.

A continuación se presentan las conclusiones específicas, ordenadas en correspondencia con los objetivos planteados y con el capítulo que desarrolla cada uno.

En relación con la fundamentación teórica (\autoref{cap:marco_teorico}), se sistematizó el MEF para la elasticidad lineal plana con elementos isoparamétricos Q4 y Q9 —de la formulación variacional a la recuperación de tensiones, con la calidad de malla y los criterios de verificación y validación— anclado a la literatura clásica de la disciplina \autocite{zienkiewicz2013fem,bathe2014fem,hughes2000fem,cook2002concepts}. Ese marco no es solo exposición: las mismas expresiones se implementan en el motor, se exhiben en los módulos educativos y se sustituyen numéricamente en la memoria de cálculo, de modo que lo que el estudiante lee, ve y calcula es una única formulación, con la convención del Jacobiano y de la extrapolación de tensiones que el código efectivamente usa.

En relación con la implementación del motor de cálculo (\autoref{cap:diseno}; verificación en el \autoref{cap:resultados}), se construyó un núcleo numérico completo en \texttt{NumPy} y \texttt{SciPy} que cubre la totalidad del canal de cálculo del MEF para elementos isoparamétricos: funciones de forma bilineales (Q4) y bicuadráticas (Q9), evaluación del Jacobiano, ensamblaje de la matriz de deformación-desplazamiento \(\bm{B}\), construcción de la matriz constitutiva \(\bm{D}\) para ambos estados planos, integración de la matriz de rigidez elemental por cuadratura de Gauss, ensamblaje global, imposición de restricciones (incluyendo desplazamientos prescritos no homogéneos) y recuperación de tensiones por extrapolación desde los puntos de Gauss \autocite{zienkiewicz2013fem,bathe2014fem,cook2002concepts}. La corrección de este motor quedó verificada mediante el método de soluciones manufacturadas (MMS) en tensión plana y en deformación plana, sobre mallas regulares y distorsionadas: las tasas de convergencia asintóticas observadas coinciden con las predicciones teóricas \(\mathcal{O}(h^{p+1})\) en norma \(L^2\) y \(\mathcal{O}(h^{p})\) en seminorma \(H^1\), esto es, \(\mathcal{O}(h^{2{,}00})\) y \(\mathcal{O}(h^{1{,}00})\) para Q4, y \(\mathcal{O}(h^{3{,}00})\) y \(\mathcal{O}(h^{2{,}00})\) para Q9, y el campo de tensiones recuperado converge a \(\mathcal{O}(h^{1{,}5})\) en Q4 y \(\mathcal{O}(h^{2})\) en Q9, con una matriz de extrapolación que reproduce exactamente los campos polinómicos de cada elemento. La obtención de estas tasas constituye evidencia rigurosa de que el motor está libre de errores de orden en la formulación de los elementos y en la recuperación de tensiones.

Respecto al diseño de la interfaz gráfica (\autoref{cap:diseno}), se logró una aplicación organizada en las tres fases canónicas del análisis —pre-proceso, proceso y post-proceso— que comparten un único lienzo interactivo. Esta decisión arquitectónica permite que el estudiante construya la geometría, asigne materiales, cargas y restricciones, resuelva el sistema y visualice los campos de desplazamiento y tensión sin cambiar de contexto visual, reforzando la continuidad conceptual entre el modelo y sus resultados. La interfaz soporta selección bidireccional entre las tablas de datos y el lienzo, deshacer y rehacer sobre el modelo completo, y validación automática de la salud del modelo previa a la resolución, todo ello con una orientación deliberadamente minimalista coherente con su propósito didáctico \autocite{lee2015interactive,bishay2020teaching}.

En cuanto a los módulos educativos (\autoref{cap:diseno}; cobertura en el \autoref{cap:resultados}), se desarrollaron componentes que exponen, paso a paso y sobre el modelo real que el estudiante ha construido, cada etapa del canal de cálculo del MEF, del mapeo isoparamétrico al ensamblaje global (\autoref{tab:modulos}). A diferencia de las visualizaciones genéricas de los textos, estos módulos operan directamente sobre el elemento seleccionado en el lienzo, conectando la abstracción matemática con la geometría concreta del problema. Este enfoque materializa el aporte pedagógico central del trabajo: hacer visible y manipulable lo que en la enseñanza tradicional permanece como una caja negra entre los datos de entrada y los resultados \autocite{lee2015eigenmodes,perezsantiago2023fem}.

La verificación y validación del software (\autoref{cap:resultados}) se abordó mediante una batería de tres niveles. La verificación de código por MMS, ya mencionada, confirmó las tasas de convergencia teóricas. La validación analítica se realizó con la viga de Timoshenko, contrastando las tensiones normales y la deflexión máxima contra la solución elástica de Timoshenko y Goodier y contra un modelo equivalente en el software comercial SAP2000 \autocite{timoshenko1970elasticity,csi2017sap2000}: frente a la solución analítica los errores fueron del \(0{,}04\,\%\) en la tensión normal \(\sigma_x\), del \(0{,}26\,\%\) en la flecha y de hasta el \(2{,}9\,\%\) en la tensión cortante; frente a SAP2000, del \(0{,}21\,\%\) en \(\sigma_x\) y de hasta el \(0{,}56\,\%\) en desplazamientos. Finalmente, la validación con la membrana de Cook \autocite{cook1974membrane} reprodujo el caso de referencia histórico y permitió ilustrar empíricamente el bloqueo por cortante (shear-locking) del elemento Q4 frente a la convergencia rápida del Q9: para igual número de grados de libertad (GDL), el Q9 alcanzó un error de \(-0{,}044\,\%\) donde el Q4 conservaba aún un \(-0{,}594\,\%\). El marco metodológico de verificación y validación adoptado sigue las prácticas establecidas en la literatura de computación científica \autocite{roache1998verification,salari2000mms}.

Por último, en cuanto a la documentación automática e interoperabilidad (\autoref{cap:diseno}; prueba de interoperabilidad en el \autoref{cap:resultados}), EduFEM genera una memoria de cálculo en PDF que reconstruye el procedimiento numérico paso a paso —desde las funciones de forma hasta las tensiones nodales— en dos estilos (educativo y directo), y ofrece el intercambio de datos mediante importación de geometría DXF y modelos en formato CSV, verificado por ida y vuelta y por idempotencia de la importación (\autoref{sec:consistencia-interna}). Esta capacidad cierra el ciclo de uso de la herramienta: el estudiante no solo experimenta con el modelo, sino que también documenta y reproduce sus cálculos.

En conjunto, los resultados confirman que EduFEM satisface los seis objetivos específicos y, con ello, el objetivo general. La herramienta combina rigor numérico verificado, una interfaz interactiva de bajo umbral de entrada y una capa pedagógica que expone el interior del método, y hace observable y contrastable cada etapa del procedimiento, que es la respuesta al problema científico planteado; su valor como apoyo a la enseñanza y el aprendizaje del MEF se sostiene por diseño y en la literatura. Confirman también la hipótesis de diseño que orientó el trabajo: fue factible construir un software que recorre el canal de cálculo completo del MEF en elasticidad plana de forma transparente, interactiva, numéricamente verificable frente a referencias de exactitud conocida y coherente con el flujo de trabajo profesional. La verificación de esa hipótesis se estableció por diseño —comprobando que la herramienta reúne dichos atributos—; la medición empírica de su impacto en el aprendizaje, frente a otras formas de enseñanza, queda planteada como trabajo futuro y se discute entre las limitaciones.

\section*{Aportes del trabajo}

El principal aporte de este trabajo es una herramienta de software libre, funcional y verificada, que articula tres capas habitualmente disociadas en los recursos educativos sobre el MEF: un motor de cálculo riguroso, una interfaz de modelado interactiva y un conjunto de módulos que explican el método sobre el modelo real del estudiante.

Desde el punto de vista numérico, el motor implementa de manera completa y verificada la formulación isoparamétrica Q4 y Q9 para los dos estados planos de la elasticidad, con un esquema de indexación de GDL robusto frente a identificadores de nodo no contiguos —que desacopla la identidad del nodo de su índice de ensamblaje— y con tratamiento de cargas volumétricas, desplazamientos prescritos no homogéneos y normas de error \(L^2\) y \(H^1\) integradas con cuadratura de orden adecuado. La validación acota el error en las magnitudes primarias por debajo del $0{,}3\,\%$ frente a la solución analítica y del $0{,}6\,\%$ frente al modelo equivalente en SAP2000, para los problemas de referencia estudiados.

Desde el punto de vista pedagógico, el aporte distintivo es la decisión de diseño de que cada módulo educativo opere sobre el elemento que el estudiante selecciona en el lienzo compartido, en lugar de sobre un ejemplo prefabricado. Esto convierte la exploración del Jacobiano, de la matriz \(\bm{B}\) o de la cuadratura de Gauss en una experiencia ligada al problema concreto, y permite que fenómenos sutiles como el bloqueo por cortante o los modos espurios (hourglass) se observen en el modelo propio. La memoria de cálculo en PDF eleva esta capa a un segundo aporte pedagógico: genera, sobre el modelo concreto del usuario, un documento que reconstruye el procedimiento completo con sustitución numérica paso a paso —de las funciones de forma a las tensiones nodales—. Frente a las visualizaciones interactivas, que muestran \emph{cómo} se calcula, la memoria deja un registro escrito que el alumno puede seguir y reproducir por sus propios medios, cerrando el lazo entre la teoría, la exploración y el cálculo manual.

\section*{Limitaciones}

El alcance del trabajo impone limitaciones que conviene explicitar. En primer lugar, el análisis se restringe a la elasticidad lineal estática en dos dimensiones (tensión y deformación plana); no se contemplan comportamientos no lineales (material o geométrico), análisis dinámico ni problemas tridimensionales; en consecuencia, el análisis estructural abordado se circunscribe a la elasticidad lineal plana de medios continuos y no alcanza a los sistemas estructurales discretos, las placas ni las cáscaras. En segundo lugar, la biblioteca de elementos se limita a los cuadriláteros isoparamétricos Q4 y Q9, sin elementos triangulares ni de orden superior, y sin mallado automático: el estudiante construye la malla manualmente, mediante importación DXF o por generación estructurada en los casos de validación.

En tercer lugar, el solucionador emplea una factorización directa sobre la matriz de rigidez almacenada en formato disperso, esquema eficiente para los tamaños de modelo propios de un contexto educativo --e incluso para mallas de decenas de miles de grados de libertad, según cuantifica el \autoref{cap:resultados}-- pero cuyo coste crece de forma desfavorable en el extremo de mallas masivas, donde un esquema iterativo o una factorización especializada resultarían preferibles. En cuarto lugar, el elemento Q4 exhibe bloqueo por cortante bajo flexión dominante, fenómeno que el software ilustra deliberadamente con fines didácticos pero que no mitiga mediante técnicas correctivas. Finalmente, aunque el diseño de la interfaz y de los módulos se guió por principios pedagógicos y por la literatura sobre enseñanza del MEF, no se realizó una validación empírica del impacto educativo con estudiantes; las conclusiones sobre la utilidad didáctica son, por tanto, de carácter cualitativo y basadas en el diseño.

\section*{Recomendaciones y trabajo futuro}

A partir de las limitaciones identificadas, se plantean líneas concretas y realistas de continuación para el proyecto, cada una anclada en el capítulo del que deriva: la ampliación de la formulación (tratamiento del bloqueo, más tipos de elemento) en el \autoref{cap:marco_teorico}; la del software (solucionador, extensión a tres dimensiones) en el \autoref{cap:diseno}; y la de la evidencia (más casos de referencia y validación pedagógica) en el \autoref{cap:resultados}.

\textbf{Optimización del solucionador.} El motor ya ensambla la matriz de rigidez global vectorizada por lotes, la almacena en formato disperso comprimido (CSR) y resuelve el sistema reducido con una factorización LU dispersa directa con ordenamiento de mínimo grado sobre $\bm{A}^{\top} + \bm{A}$, que aprovecha la simetría estructural de la matriz y reduce el tiempo de factorización y solución entre 1,7 y 2,9 veces frente al ordenamiento genérico COLAMD según el tamaño del sistema (\autoref{cap:resultados}, medido con \texttt{tests/bench\_timing.py}). Durante el desarrollo se evaluó además una reordenación de Cuthill-McKee inverso previa a la factorización y se descartó: solo compensó el costo de permutar en sistemas por encima de unos ocho mil grados de libertad, con ganancias del 7 al 19\,\%, un tamaño en el que el tiempo de resolución ya no es el factor limitante. Las extensiones de mayor impacto sobre la escalabilidad consisten, por tanto, en incorporar una factorización de Cholesky y en ofrecer un solucionador iterativo precondicionado para problemas masivos. Lo primero es aplicable porque $\bm{K}_{ff}$ es simétrica y definida positiva (\autoref{sec:ensamblaje}): al descomponerla como $\bm{K}_{ff} = \bm{L}\bm{L}^{\top}$ basta calcular y almacenar un triángulo, y la definición positiva vuelve innecesario el pivoteo, de modo que el trabajo aritmético se reduce aproximadamente a la mitad del de una factorización LU general sobre la misma matriz. Estas mejoras preservarían la interfaz pública del solucionador y la versión legible del ensamblaje reservada a los módulos educativos, de modo que no obligarían a reescribir las capas de post-proceso ni dichos módulos.

\textbf{Tratamiento del bloqueo volumétrico.} Para abordar el bloqueo (\autoref{sec:locking}) —por cortante en flexión y volumétrico cuando \(\nu \to 0{,}5\) en deformación plana— se recomienda incorporar técnicas correctivas como la integración reducida selectiva (SRI) o la formulación \(\bar{\bm{B}}\) (B-bar) \autocite{hughes2000fem,bathe2014fem}. Más allá de su valor numérico, estas técnicas tienen un alto valor pedagógico: permitirían un módulo educativo dedicado que contraste, sobre un mismo modelo, el comportamiento bloqueado y el corregido, profundizando en la comprensión del fenómeno que la herramienta ya ilustra con el caso de referencia de Cook.

\textbf{Más tipos de elemento y extensión a tres dimensiones.} La biblioteca de elementos puede ampliarse con elementos triangulares (lineales y cuadráticos) y con cuadriláteros de transición, ofreciendo así una variedad que enriquezca el estudio comparativo de la convergencia y del comportamiento bajo distorsión de malla \autocite{onate2009structural,reddy2006introduction}. A más largo plazo, la extensión del motor a problemas tridimensionales con elementos hexaédricos (H8, H20) y tetraédricos constituiría un salto natural, aprovechando que gran parte de la formulación isoparamétrica y de la infraestructura de cuadratura es generalizable a tres dimensiones.

\textbf{Ampliación de la batería de validación.} La evidencia del \autoref{cap:resultados} descansa en tres casos de referencia y en un único caso con material de ingeniería civil, la viga de hormigón. Se recomienda incorporar al menos un caso civil en deformación plana con solución de referencia publicada —una presa de gravedad, un muro de contención o un túnel somero— y un caso con carga superficial variable y peso propio, de modo que la validación externa cubra el estado plano que la práctica civil emplea con más frecuencia y las cargas que hoy solo se verifican por consistencia interna.

\textbf{Validación pedagógica con estudiantes.} La línea de continuación más importante para consolidar la finalidad del proyecto es la validación empírica de su impacto educativo. Se recomienda diseñar un estudio con estudiantes de cursos de MEF que mida, mediante una comparación antes--después e instrumentos de percepción, en qué grado el uso de EduFEM mejora la comprensión de los conceptos centrales del método respecto de la enseñanza tradicional. La formulación concreta de ese diseño excede el alcance de este trabajo y requiere el instrumental propio de la didáctica de la ingeniería; los antecedentes revisados \autocite{perezsantiago2023fem,bishay2020teaching,lee2015interactive} indican, en cambio, qué conceptos conviene poner a prueba. Los resultados de dicho estudio permitirían afinar el diseño de los módulos educativos sobre evidencia y aportarían sustento cuantitativo a la utilidad didáctica que el presente trabajo justifica por diseño.

En conjunto, estas líneas trazan un camino de evolución que preserva la identidad educativa de EduFEM al tiempo que amplía su alcance numérico, su biblioteca de elementos y, sobre todo, su fundamentación pedagógica empírica.